\chapter{Information-theoretic tools and the divergence decomposition}
\label{chap:pre_info}

This chapter collects the information-theoretic facts used by the lower bounds of
Chapters~\ref{chap:lower_adaptive} and~\ref{chap:log_gains}: the Kullback--Leibler divergence
of two Gaussians, the Pinsker and Bretagnolle--Huber inequalities, convexity of the divergence
in the pair of measures (hence in its second argument), and, most importantly, the \emph{divergence decomposition} for
algorithm-environment sequences (the ``chain rule'' of the paper, Appendix~A): the divergence
between the laws of the histories produced by one algorithm against two stationary
environments is the expected sum of the one-step divergences along the trajectory. The
decomposition is stated for the LML framework, for arbitrary reward kernels, and the linear
Gaussian case of the paper is a corollary.

\textbf{What Mathlib/LML provides.} Mathlib defines the Kullback--Leibler divergence
\texttt{InformationTheory.klDiv μ ν} as an \texttt{ℝ≥0∞}-valued \texttt{irreducible\_def}
(\texttt{Mathlib/InformationTheory/KullbackLeibler/Basic.lean}): it equals
$\int \mathrm{llr}(\mu,\nu)\,d\mu + \nu(\Omega) - \mu(\Omega)$ when $\mu \ll \nu$ and the
log-likelihood ratio $\mathrm{llr}(\mu,\nu) = \log \frac{d\mu}{d\nu}$ is $\mu$-integrable, and
$\infty$ otherwise (\texttt{klDiv\_of\_ac\_of\_integrable}, \texttt{klDiv\_of\_not\_ac},
\texttt{klDiv\_of\_not\_integrable}). It is also the $f$-divergence
$\int \mathrm{klFun}(\frac{d\mu}{d\nu})\,d\nu$ with $\mathrm{klFun}(x) = x\log x + 1 - x$
(\texttt{klDiv\_eq\_lintegral\_klFun\_of\_ac}). The chain rule is available in the form
\texttt{klDiv\_compProd\_eq\_add}:
$\KL(\mu \otimes \kappa \,\|\, \nu \otimes \eta) = \KL(\mu\|\nu) + \KL(\mu\otimes\kappa\,\|\,\mu\otimes\eta)$,
together with \texttt{klDiv\_compProd\_left} ($\KL(\mu\otimes\kappa\,\|\,\nu\otimes\kappa) = \KL(\mu\|\nu)$),
and the data-processing inequalities \texttt{klDiv\_map\_le}, \texttt{klDiv\_trim\_le},
\texttt{klDiv\_comp\_right\_le}. The lower bound \texttt{mul\_log\_le\_klDiv} and the converse
Gibbs inequality \texttt{klDiv\_eq\_zero\_iff} are also there. Radon--Nikodym derivatives of
composition-products are in \texttt{Mathlib/Probability/Kernel/Composition/RadonNikodym.lean}
(\texttt{rnDeriv\_measure\_compProd}) and kernel derivatives \texttt{Kernel.rnDeriv} are
measurable (\texttt{Mathlib/Probability/Kernel/RadonNikodym.lean}, under the hypothesis
\texttt{MeasurableSpace.CountableOrCountablyGenerated}, which is implied by a countably
generated target space); the same file proves that $\{x : \kappa_x \ll \eta_x\}$ is
measurable (\texttt{Kernel.measurableSet\_absolutelyContinuous}, via
\texttt{Kernel.singularPart\_eq\_zero\_iff\_absolutelyContinuous}). Real Gaussians are
\texttt{ProbabilityTheory.gaussianReal μ v} with density \texttt{gaussianPDFReal}
(\texttt{Mathlib/Probability/Distributions/Gaussian/Real.lean}), with
\texttt{gaussianReal\_absolutelyContinuous}, \texttt{rnDeriv\_gaussianReal},
\texttt{integral\_id\_gaussianReal}. LML provides the trajectory measure
\texttt{Learning.trajMeasure alg env}, the step decomposition of histories
(\texttt{history\_succ}, \texttt{hasCondDistrib\_step}, \texttt{hasCondDistrib\_action}),
stationary environments \texttt{stationaryEnv} (feedback kernel
\texttt{feedback\_stationaryEnv}: $(h,a) \mapsto \nu(a)$), and uniqueness of the law of the
history (\texttt{IsAlgEnvSeqUntil.map\_history}).

\textbf{What is missing.} Mathlib has neither the divergence of two Gaussians, nor Pinsker's
inequality, nor the Bretagnolle--Huber inequality, nor convexity of $\KL$ in the pair of
measures, nor the ``integrated'' form of the chain rule
$\KL(\mu\otimes\kappa\,\|\,\mu\otimes\eta) = \int \KL(\kappa_x\|\eta_x)\,\mu(dx)$ (the Mathlib file
explicitly avoids it because $x \mapsto \KL(\kappa_x\|\eta_x)$ need not be measurable in
general; it is measurable when the target space is countably generated, which is all we need).
The integrated chain rule and the invariance of $\KL$ under measurable embeddings now live in LML
(\texttt{LeanMachineLearning/ForMathlib/InformationTheory/KullbackLeibler/ChainRule.lean}); the
policy/reward specialization of Lemma~\ref{lem:pi_kl_step} and the divergence decomposition for
algorithm-environment sequences are new to LML. Mathlib also lacks the behaviour of $\KL$ under
restrictions (Lemma~\ref{lem:kl_restrict}), the invariance under sufficient statistics
(Lemma~\ref{lem:kl_sufficient_statistic}) and the continuity of $\KL$ along a generating
sequence of $\sigma$-algebras (Lemma~\ref{lem:kl_iSup_map}), which the stopping-time and
trajectory versions of the decomposition (Section~\ref{sec:pre_info_stopping}) rely on.

\section{Basic inequalities}
\label{sec:pre_info_basic}

\begin{lemma}[Invariance under measurable embeddings]\label{lem:pi_kl_measurable_equiv}
\lean{InformationTheory.klDiv_map_measurableEmbedding, InformationTheory.klDiv_map_measurableEquiv}\leanok
Let $\mu,\nu$ be finite measures on $\Omega$ and $f : \Omega \to \Omega'$ a measurable embedding
(for example a measurable equivalence). Then $\KL(f_*\mu \,\|\, f_*\nu) = \KL(\mu\|\nu)$.
\end{lemma}

\begin{proof}
\leanok
Two applications of the data-processing inequality (Mathlib \texttt{klDiv\_map\_le}): with
$f$ we get $\le$, and with a measurable left inverse $g$ of $f$ applied to $f_*\mu, f_*\nu$ we get
$\ge$, since $g_*f_*\mu = (g \circ f)_*\mu = \mu$. If $\Omega$ is empty then $\mu = \nu = 0$ and
both sides vanish; otherwise such a $g$ is given by \texttt{MeasurableEmbedding.invFun}.
\end{proof}

\begin{lemma}[Divergence of two Gaussians with the same variance]\label{lem:kl_gaussian}
\lean{ProbabilityTheory.gaussianReal_absolutelyContinuous_gaussianReal, ProbabilityTheory.llr_gaussianReal, ProbabilityTheory.integrable_llr_gaussianReal, ProbabilityTheory.klDiv_gaussianReal, ProbabilityTheory.klDiv_gaussianReal_one}\leanok
\uses{lem:gauss_1d_moments}
Let $m_1, m_2 \in \R$ and $v > 0$. Then $\Normal(m_1, v) \ll \Normal(m_2,v)$, the log-likelihood
ratio is $\mathrm{llr}(y) = \frac{(m_1 - m_2)}{v}\big(y - \frac{m_1+m_2}{2}\big)$ for Lebesgue-a.e.\
$y$, it is integrable, and
\[
  \KL\big(\Normal(m_1,v)\,\big\|\,\Normal(m_2,v)\big) = \frac{(m_1-m_2)^2}{2v} < \infty .
\]
In particular $\KL(\Normal(m_1,1)\|\Normal(m_2,1)) = (m_1-m_2)^2/2$.
\end{lemma}

\begin{proof}
\leanok
\uses{lem:gauss_1d_moments}
Both measures are $\mathrm{Leb}$ with the densities
$p_i(y) = (2\pi v)^{-1/2}\exp(-(y-m_i)^2/(2v))$ (Mathlib \texttt{gaussianReal\_of\_var\_ne\_zero}),
which are positive everywhere, so each is absolutely continuous with respect to the other and
$\frac{d\Normal(m_1,v)}{d\Normal(m_2,v)} = p_1/p_2$ a.e.\ (Mathlib
\texttt{Measure.rnDeriv\_withDensity} and the chain rule \texttt{Measure.rnDeriv\_mul\_rnDeriv}
for $\mathrm{Leb} \to \Normal(m_2,v) \to \Normal(m_1,v)$, or directly the uniqueness of the
Radon--Nikodym derivative, \texttt{Measure.eq\_rnDeriv}). Taking logarithms,
\[
  \log\frac{p_1(y)}{p_2(y)} = \frac{(y-m_2)^2 - (y-m_1)^2}{2v}
  = \frac{(m_1-m_2)\,(2y - m_1 - m_2)}{2v} ,
\]
an affine function of $y$, hence integrable under $\Normal(m_1,v)$ (Gaussians have a first
moment, Lemma~\ref{lem:gauss_1d_moments}; Mathlib \texttt{integral\_id\_gaussianReal}). Its
integral under $\Normal(m_1,v)$ is obtained by replacing $y$ by its mean $m_1$:
$\frac{(m_1-m_2)(2m_1 - m_1 - m_2)}{2v} = \frac{(m_1-m_2)^2}{2v}$. Conclude with Mathlib
\texttt{klDiv\_of\_ac\_of\_integrable} (both measures are probability measures so the correction
term $\nu(\Omega)-\mu(\Omega)$ vanishes).
\end{proof}

\begin{lemma}[Divergence of two Bernoulli measures]\label{lem:pi_kl_bernoulli}
\lean{InformationTheory.klBin, InformationTheory.klDiv_bernoulliMeasure, InformationTheory.ofReal_le_klDiv_bernoulliMeasure, ProbabilityTheory.lintegral_bernoulliMeasure, ProbabilityTheory.bernoulliMeasure_eq_withDensity, ProbabilityTheory.rnDeriv_bernoulliMeasure}\leanok
For $p \in [0,1]$ write $\mathrm{B}(p) := p\,\delta_1 + (1-p)\,\delta_0$ for the Bernoulli
measure on $\{0,1\}$. Let $p \in [0,1]$ and $q \in (0,1)$. Then
\[
  \KL(\mathrm{B}(p)\,\|\,\mathrm{B}(q)) = p\log\frac{p}{q} + (1-p)\log\frac{1-p}{1-q}
  =: \mathrm{kl}(p,q),
\]
with the convention $0 \log 0 = 0$. If $q \in \{0,1\}$ and $p \ne q$ then
$\KL(\mathrm{B}(p)\|\mathrm{B}(q)) = \infty$.
\end{lemma}

\begin{proof}
\leanok
For $q \in (0,1)$, $\mathrm{B}(q)$ charges both points so $\mathrm{B}(p) \ll \mathrm{B}(q)$, with
Radon--Nikodym derivative $p/q$ at $1$ and $(1-p)/(1-q)$ at $0$; every function on a finite set is
integrable. By Mathlib \texttt{klDiv\_eq\_lintegral\_klFun\_of\_ac}, the divergence is
$q\,\mathrm{klFun}(p/q) + (1-q)\,\mathrm{klFun}((1-p)/(1-q))$ with
$\mathrm{klFun}(x) = x\log x + 1 - x$, which expands to
$p\log(p/q) + (1-p)\log((1-p)/(1-q)) + (q - p) + ((1-q) - (1-p)) = \mathrm{kl}(p,q)$; the
convention $0\log 0 = 0$ is exactly $\mathrm{klFun}(0) = 1$ (Mathlib \texttt{klFun\_zero}). If
$q = 0$ and $p > 0$ then $\mathrm{B}(p)(\{1\}) > 0 = \mathrm{B}(0)(\{1\})$ so $\mathrm{B}(p)
\not\ll \mathrm{B}(q)$ and the divergence is $\infty$ (\texttt{klDiv\_of\_not\_ac}); the case
$q = 1$ is symmetric.
\end{proof}

\begin{lemma}[Pinsker's inequality for Bernoulli measures]\label{lem:kl_bernoulli_pinsker}
\lean{InformationTheory.klBin_pinsker}\leanok
\uses{lem:pi_kl_bernoulli}
For $p \in [0,1]$ and $q \in (0,1)$, $\mathrm{kl}(p,q) \ge 2(p-q)^2$.
\end{lemma}

\begin{proof}
\leanok
\uses{lem:pi_kl_bernoulli}
Fix $p \in (0,1)$ first and let $g(q) := \mathrm{kl}(p,q) - 2(p-q)^2$ on $(0,1)$. Then
$g(p) = 0$ and
\[
  g'(q) = -\frac{p}{q} + \frac{1-p}{1-q} + 4(p-q) = \frac{q-p}{q(1-q)} - 4(q-p)
  = (q-p)\Big(\frac{1}{q(1-q)} - 4\Big).
\]
Since $q(1-q) \le 1/4$, the last factor is nonnegative, so $g' \le 0$ on $(0,p]$ and $g' \ge 0$
on $[p,1)$: $g$ is nonincreasing then nondecreasing (Mathlib
\texttt{antitoneOn\_of\_deriv\_nonpos}, \texttt{monotoneOn\_of\_deriv\_nonneg}) and its minimum
is $g(p) = 0$. For $p = 1$ the claim is $-\log q \ge 2(1-q)^2$: the function
$h(q) := -\log q - 2(1-q)^2$ satisfies $h(1) = 0$ and $h'(q) = -1/q + 4(1-q) = -(2q-1)^2/q \le 0$,
so $h \ge 0$ on $(0,1]$. The case $p = 0$ is the case $p = 1$ with $q$ replaced by $1-q$.
\end{proof}

\begin{lemma}[Pinsker's inequality]\label{lem:pinsker}
\lean{InformationTheory.pinsker_measureReal, InformationTheory.abs_sub_le_sqrt_klDiv}\leanok
\uses{lem:pi_kl_bernoulli, lem:kl_bernoulli_pinsker}
Let $\mu, \nu$ be probability measures on a measurable space and $A$ a measurable set. Then
\[
  2\,\big(\mu(A) - \nu(A)\big)^2 \le \KL(\mu\|\nu)
\]
(as elements of $[0,\infty]$). In particular, if $\KL(\mu\|\nu) < \infty$ then
$\abs{\mu(A) - \nu(A)} \le \sqrt{\KL(\mu\|\nu)/2}$, and if $\KL(\mu\|\nu) = \infty$ the
statement is empty.
\end{lemma}

\begin{proof}
\leanok
\uses{lem:pi_kl_bernoulli, lem:kl_bernoulli_pinsker}
Let $g := \indic_A : \Omega \to \{0,1\}$, measurable. Then $g_*\mu = \mathrm{B}(p)$ and
$g_*\nu = \mathrm{B}(q)$ with $p := \mu(A)$, $q := \nu(A)$, and the data-processing inequality
(Mathlib \texttt{klDiv\_map\_le}) gives $\KL(\mathrm{B}(p)\|\mathrm{B}(q)) \le \KL(\mu\|\nu)$.
If $p = q$ there is nothing to prove. If $q \in \{0,1\}$ and $p \ne q$, then
$\KL(\mathrm{B}(p)\|\mathrm{B}(q)) = \infty$ by Lemma~\ref{lem:pi_kl_bernoulli}, so
$\KL(\mu\|\nu) = \infty$ and the inequality holds trivially. Otherwise $q \in (0,1)$ and
Lemmas~\ref{lem:pi_kl_bernoulli} and~\ref{lem:kl_bernoulli_pinsker} give
$2(p-q)^2 \le \mathrm{kl}(p,q) = \KL(\mathrm{B}(p)\|\mathrm{B}(q)) \le \KL(\mu\|\nu)$.
\end{proof}

\begin{lemma}[Bretagnolle--Huber inequality]\label{lem:bretagnolle_huber}
\lean{InformationTheory.bretagnolle_huber}\leanok
Let $\mu, \nu$ be probability measures on a measurable space and $A$ a measurable set. Then
\[
  \mu(A) + \nu(A^c) \ge \tfrac12 \exp\big(-\KL(\mu\|\nu)\big),
\]
where the right-hand side is $0$ if $\KL(\mu\|\nu) = \infty$.
\end{lemma}

\begin{proof}
\leanok
If $\KL(\mu\|\nu) = \infty$ there is nothing to prove; assume $\mu \ll \nu$ and that
$\mathrm{llr}(\mu,\nu)$ is $\mu$-integrable, so that $\KL(\mu\|\nu) = \int \mathrm{llr}\,d\mu$
(Mathlib \texttt{klDiv\_of\_ac\_of\_integrable}). Let $\lambda := \mu + \nu$ and
$p := d\mu/d\lambda$, $q := d\nu/d\lambda$, which exist (Mathlib
\texttt{Measure.rnDeriv}, \texttt{withDensity\_rnDeriv\_eq} since $\mu,\nu \ll \lambda$) and
satisfy $p + q = 1$ $\lambda$-a.e.\ (\texttt{Measure.rnDeriv\_add}, \texttt{rnDeriv\_self}), in
particular $p, q \le 1$. The proof has three steps.
\begin{enumerate}
  \item[(i)] $\mu(A) + \nu(A^c) = \int_A p\,d\lambda + \int_{A^c} q\,d\lambda \ge \int \min(p,q)\,d\lambda$.
  \item[(ii)] By the Cauchy--Schwarz inequality in $L^2(\lambda)$ applied to
        $\sqrt{\min(p,q)}$ and $\sqrt{\max(p,q)}$, whose product is $\sqrt{pq}$,
        \[
          \Big(\int\sqrt{pq}\,d\lambda\Big)^2 \le \int\min(p,q)\,d\lambda\cdot\int\max(p,q)\,d\lambda
          \le 2\int \min(p,q)\,d\lambda ,
        \]
        since $\max(p,q) \le p + q$ and $\int (p+q)\,d\lambda = \mu(\Omega) + \nu(\Omega) = 2$.
  \item[(iii)] Since $\mu \ll \nu$, we have $p = 0$ $\lambda$-a.e.\ on $\{q = 0\}$, and
        $\frac{d\mu}{d\nu} = p/q$ $\mu$-a.e.\ (Mathlib \texttt{Measure.rnDeriv\_mul\_rnDeriv}:
        $\frac{d\mu}{d\nu}\cdot\frac{d\nu}{d\lambda} = \frac{d\mu}{d\lambda}$ $\lambda$-a.e.). Hence
        \[
          \int \sqrt{pq}\,d\lambda = \int_{\{p>0\}} \sqrt{q/p}\;p\,d\lambda
          = \int \sqrt{q/p}\,d\mu = \int \exp\big(-\tfrac12\,\mathrm{llr}(\mu,\nu)\big)\,d\mu
          \ge \exp\Big(-\tfrac12\int \mathrm{llr}(\mu,\nu)\,d\mu\Big)
          = e^{-\KL(\mu\|\nu)/2},
        \]
        by Jensen's inequality for the convex function $\exp$ and the probability measure
        $\mu$ (Mathlib \texttt{ConvexOn.map\_integral\_le}, with $\mathrm{llr}$ integrable).
\end{enumerate}
Combining, $\mu(A) + \nu(A^c) \ge \int\min(p,q)\,d\lambda \ge \tfrac12 (\int\sqrt{pq}\,d\lambda)^2
\ge \tfrac12 e^{-\KL(\mu\|\nu)}$. This is Theorem~14.2 of~\cite{lattimore2020bandit}; see also
\cite{tsybakov2009introduction}.
\end{proof}

\begin{lemma}[Convexity of the divergence in the pair of measures]\label{lem:kl_joint_convex}
\lean{InformationTheory.klDiv_finsetSum_smul_le, InformationTheory.klDiv_sum_smul_le, InformationTheory.klDiv_smul_add_smul_le}\leanok
Let $\mu_1,\dots,\mu_J$ and $\nu_1,\dots,\nu_J$ be finite measures on the same space and let
$c_1,\dots,c_J \ge 0$. Then
\[
  \KL\Big(\sum_{j=1}^J c_j\mu_j \,\Big\|\, \sum_{j=1}^J c_j\nu_j\Big)
  \le \sum_{j=1}^J c_j\,\KL(\mu_j\|\nu_j) \quad\text{in } [0,\infty].
\]
With $\sum_j c_j = 1$ this is the convexity of $\KL$ in the pair of measures; that hypothesis is
not needed, since $\KL$ is positively homogeneous.
\end{lemma}

\begin{proof}
\leanok
\uses{lem:kl_compProd_kernel}
Let $\beta := \sum_j c_j\delta_j$ be the measure with weights $c$ on the index set
$\{1,\dots,J\}$, equipped with the discrete $\sigma$-algebra, and let $\kappa,\eta$ be the kernels
from the index set with $\kappa_j = \mu_j$ and $\eta_j = \nu_j$ (any function on a countable space
with measurable singletons is a kernel, Mathlib \texttt{Kernel.ofFunOfCountable}; they are finite
kernels because $\kappa_j(\Omega) \le \sum_i \mu_i(\Omega) < \infty$). Then
$\kappa\circ\beta = \sum_j c_j\mu_j$ and $\eta\circ\beta = \sum_j c_j\nu_j$, and $\kappa\circ\beta$
is the image of $\beta\otimes\kappa$ under the second projection (Mathlib
\texttt{Measure.snd\_compProd}). The data-processing inequality (Mathlib \texttt{klDiv\_map\_le})
therefore gives
\[
  \KL(\kappa\circ\beta \,\|\, \eta\circ\beta) \le \KL(\beta\otimes\kappa \,\|\, \beta\otimes\eta) ,
\]
and, the index set being countable, Lemma~\ref{lem:kl_compProd_kernel} (stated in the next
section) identifies the right-hand side with
\[
  \int^{-} \KL(\kappa_j\|\eta_j)\,\beta(dj) = \sum_{j=1}^J c_j\,\KL(\mu_j\|\nu_j) .
\]
\end{proof}

\begin{lemma}[Convexity of the divergence in its second argument]\label{lem:kl_mixture_convex}
\lean{MeasureTheory.isProbabilityMeasure_sum_smul, InformationTheory.klDiv_finsetSum_smul_right_le, InformationTheory.klDiv_sum_smul_right_le}\leanok
Let $\mu$ be a finite measure, $\nu_1,\dots,\nu_J$ finite measures on the same space,
$c_1,\dots,c_J \ge 0$ with $\sum_j c_j = 1$, and $\nu := \sum_j c_j\nu_j$. Then
\[
  \KL(\mu\|\nu) \le \sum_{j=1}^J c_j\,\KL(\mu\|\nu_j) \quad\text{in } [0,\infty].
\]
If moreover the $\nu_j$ are probability measures, then so is $\nu$.
\end{lemma}

\begin{proof}
\leanok
\uses{lem:kl_joint_convex}
Apply Lemma~\ref{lem:kl_joint_convex} with $\mu_j := \mu$ for every $j$: the mixture on the left
is $\sum_j c_j\mu = \big(\sum_j c_j\big)\mu = \mu$. The mass of $\nu$ is $\sum_j c_j = 1$.
\end{proof}

\textbf{Lean remark.} In Mathlib, $\KL$ takes values in \texttt{ℝ≥0∞}; the statements above
are therefore stated in $[0,\infty]$ with \texttt{ENNReal.ofReal} on the real side (e.g.\
\texttt{ENNReal.ofReal (2 * (μ.real A - ν.real A)\^{}2) ≤ klDiv μ ν}), and their real versions are
obtained through \texttt{klDiv\_ne\_top\_iff} and \texttt{ENNReal.toReal}. The mixture in
Lemma~\ref{lem:kl_mixture_convex} is the measure \texttt{∑ j, (c j : ℝ≥0) • ν j}.

\section{The integrated chain rule}
\label{sec:pre_info_chain}

\begin{lemma}[Chain rule with an integrated conditional divergence]\label{lem:kl_compProd_kernel}
\lean{InformationTheory.measurable_klDiv_kernel, InformationTheory.klDiv_compProd_right_eq_lintegral, InformationTheory.klDiv_compProd_eq_add_lintegral}\leanok
Let $\mu$ be a finite measure on $\mathcal{H}$ and $\kappa,\eta$ be Markov kernels from
$\mathcal{H}$ to a countably generated measurable space $\mathcal{Y}$. Then
$x \mapsto \KL(\kappa_x\|\eta_x)$ is measurable and
\[
  \KL(\mu\otimes\kappa\,\|\,\mu\otimes\eta) = \int^{-} \KL(\kappa_x\|\eta_x)\,\mu(dx) ,
\]
where $\int^-$ denotes the Lebesgue integral of a $[0,\infty]$-valued function. Consequently,
for any finite measure $\nu$ on $\mathcal{H}$,
$\KL(\mu\otimes\kappa\,\|\,\nu\otimes\eta) = \KL(\mu\|\nu) + \int^-\KL(\kappa_x\|\eta_x)\,\mu(dx)$.
\end{lemma}

\begin{proof}
\leanok
Mathlib \texttt{Measure.absolutelyContinuous\_compProd\_right\_iff} gives
$\mu\otimes\kappa \ll \mu\otimes\eta$ iff $\kappa_x \ll \eta_x$ for $\mu$-a.e.\ $x$.

\emph{Measurability.} Mathlib's kernel Radon--Nikodym theory
(\texttt{Mathlib/Probability/Kernel/RadonNikodym.lean}) is developed under the hypothesis
\texttt{MeasurableSpace.CountableOrCountablyGenerated 𝓗 𝓨}, which is implied by
\texttt{CountablyGenerated 𝓨} (an instance). Under it, the kernel Radon--Nikodym
derivative $\frac{d\kappa}{d\eta} : \mathcal{H}\times\mathcal{Y}\to[0,\infty]$ is jointly
measurable (\texttt{Kernel.rnDeriv}, \texttt{Kernel.measurable\_rnDeriv}) and
$\frac{d\kappa}{d\eta}(x,\cdot) = \frac{d\kappa_x}{d\eta_x}$ $\eta_x$-a.e.\ for every $x$
(\texttt{Kernel.rnDeriv\_eq\_rnDeriv\_measure}). The set $\{x : \kappa_x \ll \eta_x\}$ is
measurable: this is exactly Mathlib \texttt{Kernel.measurableSet\_absolutelyContinuous} (for
finite kernels), whose proof rewrites the condition as
$\kappa.\mathrm{singularPart}(\eta)(x) = 0$ by
\texttt{Kernel.singularPart\_eq\_zero\_iff\_absolutelyContinuous} and uses the measurability
of $x \mapsto \kappa.\mathrm{singularPart}(\eta)(x)(\mathcal{Y})$. On this set,
$\KL(\kappa_x\|\eta_x) = \int^-\mathrm{klFun}\big(\tfrac{d\kappa}{d\eta}(x,y)\big)\,\eta_x(dy)$
(\texttt{klDiv\_eq\_lintegral\_klFun\_of\_ac}), which is measurable in $x$ by
\texttt{Measurable.lintegral\_kernel\_prod\_right}; off this set it is the constant $\infty$.

\emph{Equality.} If $\mu\otimes\kappa \not\ll \mu\otimes\eta$, then $\kappa_x\not\ll\eta_x$ on a
set of positive $\mu$-measure, where $\KL(\kappa_x\|\eta_x) = \infty$: both sides are $\infty$.
Otherwise, by \texttt{klDiv\_eq\_lintegral\_klFun\_of\_ac} and the chain rule for derivatives
of composition-products (\texttt{rnDeriv\_measure\_compProd}, again under
\texttt{CountableOrCountablyGenerated}:
$\frac{d(\mu\otimes\kappa)}{d(\mu\otimes\eta)}(x,y) = \frac{d\kappa}{d\eta}(x,y)$ for
$(\mu\otimes\eta)$-a.e.\ $(x,y)$),
\[
  \KL(\mu\otimes\kappa\,\|\,\mu\otimes\eta)
  = \int^- \mathrm{klFun}\Big(\frac{d\kappa}{d\eta}(x,y)\Big)\,(\mu\otimes\eta)(d(x,y))
  = \int^-\!\!\int^- \mathrm{klFun}\Big(\frac{d\kappa}{d\eta}(x,y)\Big)\,\eta_x(dy)\,\mu(dx)
\]
by Tonelli for composition-products (\texttt{Measure.lintegral\_compProd}), and the inner
integral is $\KL(\kappa_x\|\eta_x)$ for $\mu$-a.e.\ $x$ (those with $\kappa_x\ll\eta_x$). The
last claim is Mathlib \texttt{klDiv\_compProd\_eq\_add} followed by the equality just proved.
\end{proof}

\begin{lemma}[Sufficient statistics and the divergence]\label{lem:kl_sufficient_statistic}
\lean{InformationTheory.klDiv_map_of_leftInverse, InformationTheory.klDiv_map_of_eq_withDensity_comp, InformationTheory.klDiv_compProd_comap, MeasureTheory.Measure.map_compProd_comap, MeasureTheory.Measure.map_compProd_of_forall_map_eq, ProbabilityTheory.HasCondDistrib.map_of_forall_map_eq, ProbabilityTheory.HasCondDistrib.prodMk_const_left, ProbabilityTheory.HasCondDistrib.restrict_preimage, MeasureTheory.Measure.restrict_compProd_prod_univ}\leanok
\uses{lem:pi_kl_measurable_equiv}
Let $\mu,\nu$ be finite measures on $\Omega$.
\begin{enumerate}
\item[(i)] If $g : \Omega \to \Omega'$ is measurable with a measurable left inverse, then
  $\KL(g_*\mu\|g_*\nu) = \KL(\mu\|\nu)$.
\item[(ii)] If $\mu = (f\circ g)\cdot\nu$ for a measurable $g : \Omega\to\Omega'$ and a measurable
  $f : \Omega' \to [0,\infty]$ (the density factors through $g$, i.e.\ $g$ is a sufficient
  statistic), then $\KL(g_*\mu\|g_*\nu) = \KL(\mu\|\nu)$.
\item[(iii)] Let $\kappa,\eta$ be finite kernels from $\mathcal{B}$ to $\mathcal{C}$ and
  $f : \Omega\to\mathcal{B}$ measurable. Then
  $\KL(\mu\otimes(\kappa\circ f)\,\|\,\mu\otimes(\eta\circ f)) = \KL(f_*\mu\otimes\kappa\,\|\,f_*\mu\otimes\eta)$:
  the conditional divergence of two kernels which depend on the conditioning variable only through
  $f$ is the conditional divergence given $f$.
\end{enumerate}
No countable generation hypothesis is needed.
\end{lemma}

\begin{proof}
\leanok
\uses{lem:pi_kl_measurable_equiv}
(i) is two applications of the data-processing inequality (as for measurable embeddings).
(ii): $g_*\mu = f\cdot g_*\nu$ (change of variables), so both pairs are absolutely continuous
with densities $f\circ g$ and $f$, and the $\mathrm{klFun}$ integrals agree by the change of
variables $g$. (iii): if $f_*\mu\otimes\kappa\not\ll f_*\mu\otimes\eta$ then, transporting along
$(\omega,c)\mapsto(f(\omega),c)$, neither is the left-hand pair, and both divergences are
infinite. Otherwise let $D$ be the density of $f_*\mu\otimes\kappa$ with respect to
$f_*\mu\otimes\eta$; the sections of $D$ integrate to $\kappa(b,t)$ for $f_*\mu$-almost every $b$,
for each fixed measurable $t$, so $\mu\otimes(\kappa\circ f)$ and $(D\circ(f\times\mathrm{id}))\cdot(\mu\otimes(\eta\circ f))$
agree on rectangles, hence everywhere, and (ii) applies with $g = f\times\mathrm{id}$.
\end{proof}

\begin{lemma}[One step of an algorithm-environment pair]\label{lem:pi_kl_step}
\lean{ProbabilityTheory.Kernel.compProd_prodMkLeft_apply, InformationTheory.klDiv_compProd_compProd_prodMkLeft_eq_klDiv_comp_compProd, InformationTheory.klDiv_compProd_compProd_prodMkLeft}\leanok
\uses{lem:kl_compProd_kernel, lem:kl_sufficient_statistic}
Let $P$ be a probability measure on a measurable space $\mathcal{H}$ (``histories''), $\pi$ a
Markov kernel from $\mathcal{H}$ to a countably generated space $\mathcal{A}$ (``policy''),
and $\kappa,\kappa'$ Markov kernels from $\mathcal{A}$ to a countably generated space
$\mathcal{Y}$ (``reward kernels'').
Write $\tilde\kappa$ for the kernel $(h,a)\mapsto\kappa(a)$ from $\mathcal{H}\times\mathcal{A}$
to $\mathcal{Y}$ (LML/Mathlib \texttt{Kernel.prodMkLeft}), and similarly $\tilde\kappa'$. Then
\[
  \KL\big(P\otimes(\pi\otimes\tilde\kappa)\,\big\|\,P\otimes(\pi\otimes\tilde\kappa')\big)
  = \int^- \KL(\kappa(a)\|\kappa'(a))\,(\pi\circ P)(da)
  = \int^-\!\!\int^- \KL(\kappa(a)\|\kappa'(a))\,\pi(h,da)\,P(dh),
\]
where $\pi\circ P := \int \pi(h,\cdot)\,P(dh)$ is the law of the next action. In
composition-product form, without any countable generation hypothesis,
$\KL\big(P\otimes(\pi\otimes\tilde\kappa)\,\big\|\,P\otimes(\pi\otimes\tilde\kappa')\big)
  = \KL\big((\pi\circ P)\otimes\kappa\,\big\|\,(\pi\circ P)\otimes\kappa'\big)$.
\end{lemma}

\begin{proof}
\leanok
\uses{lem:kl_compProd_kernel}
Apply Lemma~\ref{lem:kl_compProd_kernel} to $P$ and the kernels $\pi\otimes\tilde\kappa$,
$\pi\otimes\tilde\kappa'$, Markov kernels from $\mathcal{H}$ to the countably generated space
$\mathcal{A}\times\mathcal{Y}$ (Mathlib \texttt{MeasurableSpace.CountablyGenerated} is stable
under products). For each $h$, the measure
$(\pi\otimes\tilde\kappa)(h)$ is the composition-product $\pi(h)\otimes\kappa$ of the measure
$\pi(h)$ on $\mathcal{A}$ with the kernel $\kappa$ (Mathlib
\texttt{Kernel.compProd\_apply} and \texttt{Kernel.prodMkLeft\_apply}). By
\texttt{klDiv\_compProd\_eq\_add}, \texttt{klDiv\_self} and Lemma~\ref{lem:kl_compProd_kernel}
again,
\[
  \KL\big(\pi(h)\otimes\kappa\,\big\|\,\pi(h)\otimes\kappa'\big)
  = \KL(\pi(h)\|\pi(h)) + \int^-\KL(\kappa(a)\|\kappa'(a))\,\pi(h,da)
  = \int^-\KL(\kappa(a)\|\kappa'(a))\,\pi(h,da).
\]
Integrating in $h$ and using $\int^-\!\int^- f(a)\,\pi(h,da)\,P(dh) = \int^- f\,d(\pi\circ P)$
(Mathlib \texttt{Measure.lintegral\_bind} / \texttt{Kernel.lintegral\_comp}) gives the claim.
\end{proof}

\section{The divergence decomposition}
\label{sec:pre_info_decomposition}

\begin{definition}[A pair of stationary environments and the laws of histories]\label{def:env_pair}
Let $\mathcal{A}$ (actions) and $\mathcal{Y}$ (observations) be countably generated measurable
spaces (for instance Borel subsets of $\R^d$, $\R$, or finite types; in particular
$\mathcal{A}$ may be any action set $\Xs$ in the sense of Definition~\ref{def:arm_set}, i.e.\
a nonempty compact subset of $\R^d$, spanning or not). Let $\kappa, \kappa'$ be Markov kernels from $\mathcal{A}$ to
$\mathcal{Y}$ and let $\nu := \mathrm{stationaryEnv}(\kappa)$, $\nu' := \mathrm{stationaryEnv}(\kappa')$
be the corresponding stationary environments (LML \texttt{Learning.stationaryEnv}: the
observation at every round has law $\kappa(x_t)$ given the past and the action $x_t$). Let
$\mathrm{alg}$ be an LML algorithm with actions in $\mathcal{A}$ and observations in $\mathcal{Y}$,
with policy kernels $\pi_n$ (from histories $(x_t,y_t)_{t< n}$ of the first $n$ rounds to
$\mathcal{A}$; the initial law is $p_0 = \pi_0(\emptyset)$). For $n \ge 0$ let
\[
  \Prob^{(n)} := \mathrm{law}\big((x_t,y_t)_{t < n}\big)
  = (\mathrm{trajMeasure}(\mathrm{alg},\nu))_*\mathrm{hist}_n,
  \qquad
  \Prob'^{(n)} := (\mathrm{trajMeasure}(\mathrm{alg},\nu'))_*\mathrm{hist}_n ,
\]
the laws of the history of the first $n$ rounds (a history of \emph{length} $n$) under the two
environments. We write $\E^{(n)}$ for the expectation under $\Prob^{(n)}$ and also
$\Prob^T := \Prob^{(T)}$ when we want to emphasize the length. By LML
\texttt{IsAlgEnvSeqUntil.map\_history}, $\Prob^{(n)}$ is the law of
$(x_t,y_t)_{t< n}$ for \emph{any} algorithm-environment sequence of $(\mathrm{alg},\nu)$ on any
probability space.
\end{definition}

\begin{theorem}[Divergence decomposition]\label{thm:divergence_decomposition}
\lean{Learning.IsAlgEnvSeq.map_history_succ, Learning.IsAlgEnvSeq.klDiv_map_history, Learning.IsAlgEnvSeq.klDiv_map_history_compProd, Learning.IsAlgEnvSeq.klDiv_map_history_stepKernel}\leanok
\uses{def:env_pair, lem:stopped_history_chain_rule}
In the setting of Definition~\ref{def:env_pair}, for every $T \ge 0$,
\[
  \KL\big(\Prob^{(T)}\,\big\|\,\Prob'^{(T)}\big)
  = \sum_{t<T} \int^- \KL\big(\kappa(x_t)\,\big\|\,\kappa'(x_t)\big)\,d\Prob^{(T)}
  \qquad\text{in } [0,\infty],
\]
where $x_t$ denotes the $t$-th action of the history, that is
$\KL(\Prob^T\|\Prob'^T) = \sum_{t<T}\E^{(T)}[\KL(\kappa(x_t)\|\kappa'(x_t))]$.
If moreover $\KL(\kappa(a)\|\kappa'(a)) \le C < \infty$ for all $a \in \mathcal{A}$, then
$\KL(\Prob^T\|\Prob'^T) \le TC < \infty$ and the identity holds in $\R$ with ordinary
expectations.
\end{theorem}

\begin{proof}
\leanok
\uses{lem:pi_kl_measurable_equiv, lem:kl_compProd_kernel, lem:pi_kl_step, lem:stopped_history_chain_rule}
In the formalization the theorem is the special case $S = \emptyset$ (no stopping) of the chain
rule for stopped histories, Lemma~\ref{lem:stopped_history_chain_rule}: the stopping time is
$+\infty$, the history stopped at $\min(\tau, T)$ is the history of the first $T$ rounds, and the
one-step terms are those of Lemma~\ref{lem:pi_kl_step}. We nevertheless record the direct
induction, which is the argument of the paper.

Write $\mathcal{H}_n := (\{0,\dots,n-1\} \to \mathcal{A}\times\mathcal{Y})$ for the space of
histories of the first $n$ rounds. By definition of the stationary environment, the feedback
kernel at every round is $\tilde\kappa : (h,a)\mapsto\kappa(a)$ (LML
\texttt{feedback\_stationaryEnv}). We prove the identity by induction on $T$.

\emph{Base case $T = 0$.} $\mathcal{H}_0$ is a one-point space, so $\Prob^{(0)} = \Prob'^{(0)}$
is the Dirac mass at the empty history (LML \texttt{hasLaw\_history\_zero}) and both sides
vanish (\texttt{klDiv\_self} and the empty sum).

\emph{Induction step.} Assume the identity for $T$. By LML \texttt{history\_succ} and
\texttt{hasCondDistrib\_step}, $\mathrm{hist}_{T+1} = e^{-1}(\mathrm{hist}_T, (x_T,y_T))$
for the measurable equivalence $e : \mathcal{H}_{T+1}\simeq\mathcal{H}_T\times(\mathcal{A}\times\mathcal{Y})$
(\texttt{MeasurableEquiv.finSuccProd}), and the law of $(\mathrm{hist}_T,(x_T,y_T))$ is
the composition-product $\Prob^{(T)}\otimes(\pi_T\otimes\tilde\kappa)$ (the step kernel
\texttt{stepKernel alg env n = alg.policy n ⊗ₖ env.feedback n}); similarly with primes, with
the \emph{same} policy kernel $\pi_T$. Hence, by Lemma~\ref{lem:pi_kl_measurable_equiv},
\texttt{klDiv\_compProd\_eq\_add} and Lemma~\ref{lem:pi_kl_step},
\begin{align*}
  \KL(\Prob^{(T+1)}\|\Prob'^{(T+1)})
  &= \KL\big(\Prob^{(T)}\otimes(\pi_T\otimes\tilde\kappa)\,\big\|\,\Prob'^{(T)}\otimes(\pi_T\otimes\tilde\kappa')\big)\\
  &= \KL(\Prob^{(T)}\|\Prob'^{(T)})
     + \KL\big(\Prob^{(T)}\otimes(\pi_T\otimes\tilde\kappa)\,\big\|\,\Prob^{(T)}\otimes(\pi_T\otimes\tilde\kappa')\big)\\
  &= \sum_{t<T}\int^-\KL(\kappa(x_t)\|\kappa'(x_t))\,d\Prob^{(T)}
     + \int^-\KL(\kappa(a)\|\kappa'(a))\,(\pi_T\circ\Prob^{(T)})(da).
\end{align*}
Finally $\pi_T\circ\Prob^{(T)}$ is the law of $x_T$ under $\Prob^{(T+1)}$ (LML
\texttt{hasCondDistrib\_action}: the conditional law of $x_T$ given $\mathrm{hist}_T$ is
$\pi_T(\mathrm{hist}_T)$), and $\Prob^{(T)}$ is the image of $\Prob^{(T+1)}$ under the restriction
$\mathcal{H}_{T+1}\to\mathcal{H}_T$ (consistency of the laws of histories, LML
\texttt{history\_eq\_comp\_history}), so every term is an integral under $\Prob^{(T+1)}$ of a
function of the history, and the sum runs over $t < T+1$.

\emph{Finite case.} If the one-step divergences are bounded by $C$, the right-hand side is at
most $(n+1)C$, so all quantities are finite and \texttt{ENNReal.toReal} turns the identity into
a real identity, the integrals becoming Bochner integrals of bounded measurable functions
(\texttt{integral\_eq\_lintegral\_of\_nonneg\_ae}).
\end{proof}

\textbf{Lean remark.} The formalization (\texttt{Learning.IsAlgEnvSeq.klDiv\_map\_history})
states the theorem for two algorithm-environment sequences \texttt{IsAlgEnvSeq X Y alg
(stationaryEnv κ) P} and \texttt{IsAlgEnvSeq X' Y' alg (stationaryEnv κ') P'} on arbitrary
probability spaces, as \texttt{klDiv (P.map (history X Y T)) (P'.map (history X' Y' T)) =
∑ t ∈ range T, ∫⁻ ω, klDiv (κ (X t ω)) (κ' (X t ω)) ∂P}; this covers the trajectory
measures (which are such sequences) and avoids the uniqueness-of-law step. The finite form is
only recorded for the Gaussian case (Corollary~\ref{cor:divergence_decomposition_gaussian}).
Every conditional divergence is recorded in two forms: the \emph{integral form} above, which
needs the hypothesis \texttt{MeasurableSpace.CountablyGenerated 𝓨} for the measurability of the
integrands (Lemma~\ref{lem:kl_compProd_kernel}; it holds for $\mathcal{Y} = \R$), and the
\emph{composition-product form} \texttt{∑ t ∈ range T, klDiv (P.map (X t) ⊗ₘ κ) (P.map
(X t) ⊗ₘ κ')} (\texttt{klDiv\_map\_history\_compProd}), valid on arbitrary measurable spaces.
The chain rule for two arbitrary environments, with the conditional divergences of the step
kernels, is \texttt{klDiv\_map\_history\_stepKernel}. Only the environment changes: the policy
kernels are shared, which is what makes \texttt{klDiv\_compProd\_left} (inside
Lemma~\ref{lem:pi_kl_step}) applicable. This is the environment analogue of LML's
\texttt{Algorithm.density}, and of Lemma~\ref{lem:env_likelihood_ratio} below, which give the
Radon--Nikodym derivative rather than its expected logarithm.

\begin{lemma}[Appending the recommendation does not change the divergence]
\label{lem:kl_data_processing_recommendation}
\lean{Learning.IdentAlg.IsFixedBudget.stoppingTime_eq, Learning.IdentAlg.IsFixedBudget.stoppedHist_eq, Learning.IdentAlg.IsFixedBudget.stoppingTime_ne_top, Learning.IdentAlg.IsRun.klDiv_map_stoppedHist_out_of_isFixedBudget, Learning.IdentAlg.IsRun.klDiv_map_out_le, Learning.LinearBandit.IsRun.klDiv_map_out_le, Learning.LinearBandit.IsRun.klDiv_map_out_le_of_sq_le, Learning.LinearBandit.IsRun.exp_neg_le_measureReal_add_of_sq_le}\leanok
\uses{def:env_pair, def:bai_alg, lem:pinsker, lem:bretagnolle_huber}
In the setting of Definition~\ref{def:env_pair}, let $T \ge 1$ and let $\rho$ be a Markov
kernel from histories of length $T$ to a measurable space $\mathcal{X}$ (a recommendation rule).
Then
\[
  \KL\big(\Prob^T\otimes\rho\,\big\|\,\Prob'^T\otimes\rho\big) = \KL(\Prob^T\|\Prob'^T),
\]
and for every measurable set $E$ of pairs (history, recommendation),
\[
  (\Prob^T\otimes\rho)(E) - (\Prob'^T\otimes\rho)(E) \le \sqrt{\KL(\Prob^T\|\Prob'^T)/2},
  \qquad
  (\Prob^T\otimes\rho)(E) + (\Prob'^T\otimes\rho)(E^c) \ge \tfrac12\exp\big(-\KL(\Prob^T\|\Prob'^T)\big).
\]
When the environments are those of Definition~\ref{def:env} on an arbitrary action set
$\Xs$ (nonempty compact, not assumed spanning) and $(\mathrm{alg},\rho)$ is an
identification algorithm (Definition~\ref{def:bai_alg}), $\Prob^T\otimes\rho$ is the law
$\Prob_\theta$ of $(\hist_T,\rec)$.
\end{lemma}

\begin{proof}
\leanok
\uses{lem:pinsker, lem:bretagnolle_huber}
The equality is Mathlib \texttt{klDiv\_compProd\_left} (same kernel on both sides). The two
inequalities are Lemma~\ref{lem:pinsker} and Lemma~\ref{lem:bretagnolle_huber} applied to the
probability measures $\Prob^T\otimes\rho$, $\Prob'^T\otimes\rho$ (trivial if the divergence is
infinite). The last sentence is the definition of $\Prob_\theta$ in
Definition~\ref{def:bai_alg} together with the uniqueness of the law of the history.
\end{proof}

\textbf{Lean remark.} In the formalization the recommendation is the output of a run
(Definition~\ref{def:bai_alg}, \texttt{IsRun}): for a fixed-budget algorithm the stopping time is
$T$ and the history at the stopping time is the history of the $T$ rounds
(\texttt{IsFixedBudget.stoppingTime\_eq}, \texttt{IsFixedBudget.stoppedHist\_eq}); the law of
(history, output) is the composition-product with the output kernel, which is a probability kernel
on the stopping set, so \texttt{klDiv\_compProd\_left} applies in the form
\texttt{klDiv\_compProd\_left\_of\_ae}. Only the inequality
$\KL(\mathrm{law}(\rec)\|\mathrm{law}'(\rec)) \le \KL(\Prob^T\|\Prob'^T)$ is recorded
(\texttt{IsRun.klDiv\_map\_out\_le}, by data processing), which is what the lower bounds use, and
the Bretagnolle--Huber consequence for linear Gaussian environments,
$\Prob_\theta(\rec\in B) + \Prob_{\theta'}(\rec\notin B) \ge \tfrac12\exp(-TR^2\norm{\theta-\theta'}^2/2)$,
is \texttt{LinearBandit.IsRun.exp\_neg\_le\_measureReal\_add}.

\subsection{Stopped histories and infinite trajectories}
\label{sec:pre_info_stopping}

The paper only needs the divergence decomposition for a fixed number of rounds. For the library
we also prove it for histories stopped at a stopping time (the change-of-measure lemma of
fixed-confidence lower bounds, Kaufmann, Capp\'e and Garivier 2016) and for whole trajectories.
The results below are formalized in \texttt{Maiti2026Power/LeanMachineLearning/\{HistoryLaw,
StoppedHistory, DivergenceDecomposition, RunDivergence\}.lean} and
\texttt{Maiti2026Power/Mathlib/InformationTheory/\{KLCompProd, KLRestrict, KLFiltration\}.lean}.

\begin{definition}[Histories of the first $n$ rounds, step kernels]\label{def:fin_history}
\lean{Learning.history, Learning.stepKernel, Learning.IsAlgEnvSeq.hasCondDistrib_action, Learning.IsAlgEnvSeq.hasCondDistrib_feedback, Learning.IsAlgEnvSeq.hasCondDistrib_step, Learning.IsAlgEnvSeq.map_history_succ, Learning.map_history_succ_of_hasCondDistrib, Learning.feedback_stationaryEnv, Learning.stepKernel_stationaryEnv}\leanok
\uses{def:env_pair}
For action and feedback processes $X, Y$ on $\Omega$, the \emph{history of the first $n$ rounds}
is $\mathrm{hist}_n := (X_t, Y_t)_{t < n} : \Omega \to (\{0,\dots,n-1\}\to\mathcal{A}\times\mathcal{Y})$
(LML \texttt{Learning.history}; $n = 0$ is allowed and gives the one-point space), and
$\mathrm{hist}_{n+1} = (\mathrm{hist}_n, (X_n, Y_n))$ (\texttt{Fin.snoc}). For an algorithm
$\mathrm{alg}$ with policy kernels $\pi_n$ and an environment $\nu$ with feedback kernels
$\nu_n$, the \emph{step kernel at round $n$} is $\sigma_n := \pi_n\otimes\nu_n$, a Markov kernel
from histories of $n$ rounds to $\mathcal{A}\times\mathcal{Y}$ (LML \texttt{Learning.stepKernel}).
For an algorithm-environment sequence, $X_n$, $Y_n$ and $(X_n, Y_n)$ have conditional laws
$\pi_n$, $\nu_n$ and $\sigma_n$ given $\mathrm{hist}_n$ (and $X_n$), and the law of
$\mathrm{hist}_{n+1}$ is $\mathrm{law}(\mathrm{hist}_n)\otimes\sigma_n$. For a stationary
environment with reward kernel $\kappa$, $\nu_n = \tilde\kappa : (h, a)\mapsto\kappa(a)$ at
every round, so $\sigma_n = \pi_n\otimes\tilde\kappa$.
\end{definition}

\begin{definition}[Stopping rules, stopping times, stopped histories]\label{def:stopping_rule}
\lean{Learning.stoppingTime, Learning.stoppedHist, Learning.stoppingTime_le_iff, Learning.lt_stoppingTime_iff, Learning.stoppingTime_eq_coe_iff, Learning.stoppingTime_eq_top_iff, Learning.stoppingTime_empty, Learning.notMem_of_lt_stoppingTime, Learning.stoppedHist_mem_of_ne_top, Learning.stoppedHist_min_of_le, Learning.stoppedHist_min_of_lt, Learning.measurable_stoppingTime, Learning.measurable_stoppedHist, Learning.IsAlgEnvSeq.isStoppingTime_stoppingTime, Learning.exists_measurableSet_preimage_lt_stoppingTime, Learning.truncHist, Learning.truncHist_stoppedHist, Learning.map_stoppedHist_min_eq_map_truncHist, Learning.restrict_map_truncHist, Learning.IdentAlg.stoppingTime, Learning.IdentAlg.stoppedHist}\leanok
\uses{def:fin_history, def:bai_alg}
A \emph{stopping rule} is a measurable set $S$ of histories of variable length,
$S \subseteq \bigsqcup_n (\{0,\dots,n-1\}\to\mathcal{A}\times\mathcal{Y})$: the interaction stops after
$n$ rounds if the history of these $n$ rounds belongs to $S$. Its \emph{stopping time}
$\tau_S := \inf\{n : \mathrm{hist}_n \in S\} \in \N\cup\{\infty\}$ (Mathlib \texttt{hittingAfter})
is the number of rounds played; it is a stopping time of the history filtration of any
algorithm-environment sequence, $\{n < \tau_S\}$ is determined by $\mathrm{hist}_n$, and
$\tau_S = \infty$ for $S = \emptyset$. For a random time $\tau$, the \emph{stopped history}
$\mathrm{hist}_\tau := \langle\tau, \mathrm{hist}_\tau\rangle$ is the history of the first $\tau$
rounds as a history of variable length (of length $0$ if $\tau = \infty$); it belongs to $S$ when
$\tau = \tau_S < \infty$. The \emph{truncation} $\mathrm{trunc}_M$ keeps the first $M$ rounds of a
history of variable length: $\mathrm{trunc}_M(\mathrm{hist}_\tau) = \mathrm{hist}_{\min(\tau,M)}$
when $\tau < \infty$, so that the law of $\mathrm{hist}_{\min(\tau,M)}$ is the image of the law of
$\mathrm{hist}_\tau$ by $\mathrm{trunc}_M$ when $\tau$ is almost surely finite, and both laws
agree on histories of length $< M$. The stopping time and stopped history of an identification
algorithm (Definition~\ref{def:bai_alg}) are those of its stopping rule.
\end{definition}

\begin{lemma}[Restrictions and the divergence]\label{lem:kl_restrict}
\lean{MeasureTheory.Measure.rnDeriv_restrict_restrict, InformationTheory.klDiv_restrict_of_ac, InformationTheory.klDiv_restrict_add_restrict_compl, InformationTheory.klDiv_restrict_le, InformationTheory.klDiv_add_add_of_measure_eq_zero, InformationTheory.klDiv_eq_iSup_restrict}\leanok
Let $\mu,\nu$ be finite measures on $\Omega$ and $s$ a measurable set. Then
$\KL(\mu\|\nu) = \KL(\mu|_s\,\|\,\nu|_s) + \KL(\mu|_{s^c}\,\|\,\nu|_{s^c})$; in particular
$\KL(\mu|_s\|\nu|_s) \le \KL(\mu\|\nu)$, and the divergence of two sums of measures carried by
$s$ and $s^c$ respectively is the sum of the divergences of the two parts. If $s_n \uparrow \Omega$
is an increasing sequence of measurable sets, $\KL(\mu\|\nu) = \sup_n \KL(\mu|_{s_n}\|\nu|_{s_n})$.
\end{lemma}

\begin{proof}
\leanok
The Radon--Nikodym derivative of $\mu|_s$ with respect to $\nu|_s$ is that of $\mu$ with respect
to $\nu$ (both are $\nu|_s$-densities of $\mu|_s = (\nu|_s)\cdot\frac{d\mu}{d\nu}$), so when
$\mu\ll\nu$ the identities are those of the integral $\int\mathrm{klFun}(d\mu/d\nu)\,d\nu$ over
$s$, $s^c$ and $s_n$ (monotone convergence for the last one, $\mathrm{klFun}\ge 0$). When
$\mu\not\ll\nu$, a set $A$ with $\nu(A) = 0 < \mu(A)$ meets $s$ or $s^c$, and some $s_n$, with
positive $\mu$-measure, so the corresponding restricted divergences are infinite too.
\end{proof}

\begin{lemma}[Chain rule for stopped histories]\label{lem:stopped_history_chain_rule}
\lean{Learning.map_stoppedHist_min_eq_add, Learning.map_stoppedHist_min_succ_eq_add, Learning.map_stoppedHist_min_zero, Learning.IsAlgEnvSeq.map_history_succ_restrict_lt_stoppingTime, Learning.IsAlgEnvSeq.map_action_restrict_lt_stoppingTime, Learning.IsAlgEnvSeq.klDiv_map_stoppedHist_min_stepKernel, Learning.IsAlgEnvSeq.klDiv_map_stoppedHist_min_compProd, Learning.IsAlgEnvSeq.klDiv_map_stoppedHist_min, Learning.sum_ite_lt_eq_sum_range_toNat_min}\leanok
\uses{def:env_pair, def:fin_history, def:stopping_rule, lem:kl_restrict, lem:pi_kl_measurable_equiv, lem:pi_kl_step, lem:kl_sufficient_statistic}
Let $\mathrm{alg}$ be an algorithm, $\nu,\nu'$ two environments, $(X,Y)$ and $(X',Y')$
algorithm-environment sequences of $(\mathrm{alg},\nu)$ and $(\mathrm{alg},\nu')$ on probability
spaces $(\Omega,\Prob)$ and $(\Omega',\Prob')$, $S$ a stopping rule with stopping times $\tau$,
$\tau'$ on the two spaces, and $M \ge 0$. Then
\[
  \KL\big(\mathrm{law}(\mathrm{hist}_{\min(\tau,M)})\,\big\|\,\mathrm{law}(\mathrm{hist}'_{\min(\tau',M)})\big)
  = \sum_{t < M} \KL\big(\Prob|_{\{t<\tau\}}\circ\mathrm{hist}_t^{-1}\otimes\sigma_t\,\big\|\,
      \Prob|_{\{t<\tau\}}\circ\mathrm{hist}_t^{-1}\otimes\sigma'_t\big),
\]
the sum over the rounds of the conditional divergences of the step kernels $\sigma_t,\sigma'_t$
of the two environments given the first $t$ rounds, on the event $\{t < \tau\}$ (the round is
played). For two stationary environments with reward kernels $\kappa,\kappa'$ the $t$-th term is
$\KL(\Prob|_{\{t<\tau\}}\circ X_t^{-1}\otimes\kappa\,\|\,\Prob|_{\{t<\tau\}}\circ X_t^{-1}\otimes\kappa')$
(composition-product form) and, when $\mathcal{Y}$ is countably generated, the right-hand side is
$\E\big[\sum_{t<\min(\tau,M)}\KL(\kappa(X_t)\|\kappa'(X_t))\big]$ (integral form).
\end{lemma}

\begin{proof}
\leanok
\uses{lem:kl_restrict, lem:pi_kl_measurable_equiv, lem:pi_kl_step, lem:kl_sufficient_statistic, lem:kl_compProd_kernel}
Induction on $M$. For $M = 0$ both laws are the Dirac mass at the empty history. For the step
$M\to M+1$, write $Z_M := \mathrm{hist}_{\min(\tau,M)}$. The law of $Z_{M+1}$ splits according
to whether $\tau \le M$: on $\{\tau\le M\}$, $Z_{M+1} = Z_M$, and on $\{M < \tau\}$, $Z_{M+1}$ is the
history of the first $M+1$ rounds; the two parts are carried by the disjoint sets of histories of
length $\le M$ and of length $M+1$. Likewise the law of $Z_M$ splits into its restriction to
$\{\tau\le M\}$ (carried by $S$, as $Z_M = \mathrm{hist}_\tau \in S$ there) and the law of
$\mathrm{hist}_M$ on $\{M<\tau\}$ (carried by $S^c$). By the additivity of
Lemma~\ref{lem:kl_restrict} over disjoint supports, the invariance under the embeddings
$h\mapsto\langle n,h\rangle$ (Lemma~\ref{lem:pi_kl_measurable_equiv}) and the chain rule
\texttt{klDiv\_compProd\_eq\_add}, since on $\{M<\tau\}$ (an event determined by
$\mathrm{hist}_M$) the law of $\mathrm{hist}_{M+1}$ is $\mathrm{law}(\mathrm{hist}_M)\otimes\sigma_M$,
\[
  \KL(Z_{M+1}) = \KL(Z_M|_{\tau\le M}) + \KL(\rho_M) + \KL(\rho_M\otimes\sigma_M\,\|\,\rho_M\otimes\sigma'_M)
  = \KL(Z_M) + \KL(\rho_M\otimes\sigma_M\,\|\,\rho_M\otimes\sigma'_M),
\]
where $\rho_M$ is the law of $\mathrm{hist}_M$ on $\{M<\tau\}$ (primes omitted on the second
arguments). The stationary case is Lemma~\ref{lem:pi_kl_step} in composition-product form, the
law of $X_t$ on $\{t<\tau\}$ being the policy applied to the law of $\mathrm{hist}_t$ on that
event; the integral form is Lemma~\ref{lem:kl_compProd_kernel} and the exchange of the finite sum
with the integral.
\end{proof}

\begin{theorem}[Divergence decomposition at a stopping time]\label{thm:divergence_decomposition_stopping}
\lean{Learning.IsAlgEnvSeq.klDiv_map_stoppedHist_stepKernel, Learning.IsAlgEnvSeq.klDiv_map_stoppedHist_compProd, Learning.IsAlgEnvSeq.klDiv_map_stoppedHist}\leanok
\uses{def:env_pair, def:stopping_rule, lem:stopped_history_chain_rule, lem:kl_restrict}
In the setting of Lemma~\ref{lem:stopped_history_chain_rule}, assume that the stopping time is
almost surely finite under both $\Prob$ and $\Prob'$. Then
\[
  \KL\big(\mathrm{law}(\mathrm{hist}_\tau)\,\big\|\,\mathrm{law}(\mathrm{hist}'_{\tau'})\big)
  = \sum_{t=0}^{\infty} \KL\big(\Prob|_{\{t<\tau\}}\circ\mathrm{hist}_t^{-1}\otimes\sigma_t\,\big\|\,
      \Prob|_{\{t<\tau\}}\circ\mathrm{hist}_t^{-1}\otimes\sigma'_t\big),
\]
and for two stationary environments with reward kernels $\kappa,\kappa'$,
\[
  \KL\big(\mathrm{law}(\mathrm{hist}_\tau)\,\big\|\,\mathrm{law}(\mathrm{hist}'_{\tau'})\big)
  = \sum_{t=0}^{\infty}\KL\big(\Prob|_{\{t<\tau\}}\circ X_t^{-1}\otimes\kappa\,\big\|\,\Prob|_{\{t<\tau\}}\circ X_t^{-1}\otimes\kappa'\big)
  = \E\Big[\sum_{t<\tau}\KL(\kappa(X_t)\|\kappa'(X_t))\Big],
\]
the last form when $\mathcal{Y}$ is countably generated.
\end{theorem}

\begin{proof}
\leanok
\uses{lem:stopped_history_chain_rule, lem:kl_restrict}
The series is the supremum over $M$ of the partial sums, which are the divergences of the laws
of $\mathrm{hist}_{\min(\tau,M)}$ by Lemma~\ref{lem:stopped_history_chain_rule}. Since the
stopping times are almost surely finite, these laws are the images of the laws
$\mu,\mu'$ of the stopped histories by the truncation $\mathrm{trunc}_M$, so their divergences
are at most $\KL(\mu\|\mu')$ (data processing). Conversely, the images agree with $\mu,\mu'$ on
histories of length $< M$, and these sets increase to the whole space, so by
Lemma~\ref{lem:kl_restrict} (monotone convergence, then monotonicity under restriction)
$\KL(\mu\|\mu') = \sup_M\KL(\mu|_{\mathrm{length}<M}\|\mu'|_{\mathrm{length}<M}) \le \sup_M\KL(\mathrm{trunc}_M{}_*\mu\|\mathrm{trunc}_M{}_*\mu')$.
The integral form follows from the composition-product form as in
Lemma~\ref{lem:stopped_history_chain_rule}, the series and the integral being exchanged by
monotone convergence, and $\sum_t\mathbf 1_{t<\tau}(\cdot) = \sum_{t<\tau}(\cdot)$ on $\{\tau<\infty\}$.
\end{proof}

\begin{lemma}[Change of measure at a stopping time for runs]\label{lem:run_change_of_measure}
\lean{Learning.IdentAlg.IsRun.ae_stoppedHist_mem_stopSet, Learning.IdentAlg.IsRun.klDiv_map_stoppedHist_out, Learning.IdentAlg.IsRun.klDiv_map_out_le_stoppedHist, Learning.IdentAlg.IsRun.klDiv_map_out_le_tsum_compProd, Learning.IdentAlg.IsRun.klDiv_map_out_le_lintegral}\leanok
\uses{def:bai_alg, def:stopping_rule, thm:divergence_decomposition_stopping, lem:kl_data_processing_recommendation}
Let $A$ be an identification algorithm and consider two runs of $A$, in two environments, on
probability spaces $(\Omega,\Prob)$ and $(\Omega',\Prob')$, whose stopping times are almost surely
finite. Then the divergence between the laws of (stopped history, output) is the divergence
between the laws of the stopped histories, and the divergence between the laws of the outputs
is at most the latter. For two stationary environments with reward kernels $\kappa,\kappa'$,
\[
  \KL\big(\mathrm{law}(\rec)\,\big\|\,\mathrm{law}'(\rec')\big)
  \le \E\Big[\sum_{t<\tau}\KL(\kappa(X_t)\|\kappa'(X_t))\Big]
\]
(when $\mathcal{Y}$ is countably generated; the composition-product form holds in general).
\end{lemma}

\begin{proof}
\leanok
\uses{thm:divergence_decomposition_stopping, lem:kl_data_processing_recommendation}
As in Lemma~\ref{lem:kl_data_processing_recommendation}: the stopped history belongs to the
stopping set almost surely, on which the output rule is a probability kernel, so
\texttt{klDiv\_compProd\_left\_of\_ae} applies; then data processing along the projection on the
output, and Theorem~\ref{thm:divergence_decomposition_stopping}.
\end{proof}

\begin{lemma}[The divergence along a generating sequence of maps]\label{lem:kl_iSup_map}
\lean{InformationTheory.klDiv_map_eq_lintegral_klFun_condExp, InformationTheory.klDiv_eq_iSup_map, MeasurableSpace.iSup_comap_restrictFin}\leanok
\uses{lem:kl_restrict}
Let $\mu,\nu$ be finite measures on $\Omega$ and $g_n : \Omega\to\Omega_n$ be measurable maps whose
$\sigma$-algebras $\sigma(g_n)$ increase to the $\sigma$-algebra of $\Omega$. Then
$\KL(\mu\|\nu) = \sup_n \KL({g_n}_*\mu\,\|\,{g_n}_*\nu)$. The projections of a sequence space
$\N\to E$ on the coordinates $\le n$ satisfy the hypothesis.
\end{lemma}

\begin{proof}
\leanok
\uses{lem:kl_restrict}
The inequality $\ge$ is the data-processing inequality. For $\le$, if $\mu\ll\nu$ with density
$f$, then $\KL({g_n}_*\mu\|{g_n}_*\nu) = \int\mathrm{klFun}(\E_\nu[f\mid\sigma(g_n)])\,d\nu$ (Mathlib
\texttt{toReal\_rnDeriv\_map}), the conditional expectations converge $\nu$-almost everywhere to
$f$ by L\'evy's upward theorem (Mathlib \texttt{Integrable.tendsto\_ae\_condExp}), and Fatou's
lemma gives $\int\mathrm{klFun}(f)\,d\nu \le \liminf_n\int\mathrm{klFun}(\E_\nu[f\mid\sigma(g_n)])\,d\nu$.
If $\mu\not\ll\nu$, pick $A$ with $\nu(A) = 0 < \mu(A)$; by L\'evy's theorem under $\mu+\nu$
applied to $\mathbf 1_A$, the $\sigma(g_n)$-measurable sets $B_n := \{\E_{\mu+\nu}[\mathbf 1_A\mid\sigma(g_n)] > 1/2\}$
satisfy $\mu(B_n)\to\mu(A)$ and $\nu(B_n)\to 0$, and the lower bound
$\mu(B)\log(\mu(B)/\nu(B)) + \nu(B) - \mu(B) \le \KL(\mu|_B\|\nu|_B) \le \KL(\mu\|\nu)$ (Mathlib
\texttt{mul\_log\_le\_klDiv} and Lemma~\ref{lem:kl_restrict}), applied to the images and the
sets $B_n = g_n^{-1}(C_n)$, shows that the divergences of the images tend to infinity.
\end{proof}

\begin{theorem}[Divergence decomposition for trajectories]\label{thm:divergence_decomposition_trajectory}
\lean{Learning.klDiv_map_trajectory_eq_iSup, Learning.IsAlgEnvSeq.klDiv_map_trajectory_stepKernel, Learning.IsAlgEnvSeq.klDiv_map_trajectory_compProd, Learning.IsAlgEnvSeq.klDiv_map_trajectory}\leanok
\uses{def:env_pair, thm:divergence_decomposition, lem:kl_iSup_map}
In the setting of Definition~\ref{def:env_pair}, for two algorithm-environment sequences on
$(\Omega,\Prob)$, $(\Omega',\Prob')$, the divergence between the laws of the whole trajectories
$(X_t,Y_t)_{t\in\N}$ is
\[
  \sum_{t=0}^{\infty}\KL\big(\Prob\circ X_t^{-1}\otimes\kappa\,\big\|\,\Prob\circ X_t^{-1}\otimes\kappa'\big)
  = \sum_{t=0}^{\infty}\E\big[\KL(\kappa(X_t)\|\kappa'(X_t))\big]
\]
(the last form when $\mathcal{Y}$ is countably generated), and for two arbitrary environments it
is the series of the conditional divergences of the step kernels given the past.
\end{theorem}

\begin{proof}
\leanok
\uses{thm:divergence_decomposition, lem:kl_iSup_map}
By Lemma~\ref{lem:kl_iSup_map} applied to the projections on the first $n+1$ coordinates, the
divergence between the laws of the trajectories is the supremum over $n$ of the divergences
between the laws of the histories up to time $n$, which are the partial sums of the series by
Theorem~\ref{thm:divergence_decomposition}.
\end{proof}

\begin{corollary}[Divergence decomposition for linear Gaussian environments]
\label{cor:divergence_decomposition_gaussian}
\lean{Learning.LinearBandit.klDiv_linearGaussianKernel, Learning.LinearBandit.klDiv_map_history, Learning.LinearBandit.klDiv_map_history_le}\leanok
\uses{def:env, def:env_pair, thm:divergence_decomposition, lem:kl_gaussian}
Let $\Xs$ be an action set (Definition~\ref{def:arm_set}: nonempty compact, \emph{not}
assumed spanning; Part~\ref{part:main} applies this corollary to non-spanning sets) and
$\theta,\theta'\in\R^d$. For every
algorithm with actions in $\Xs$ and every $T \ge 1$, writing $\Prob_\theta^T$ for the law of the
history of length $T$ under $\nu_\theta$ (Definition~\ref{def:env}; this is $\Prob^T$ of
Definition~\ref{def:env_pair} with $\mathcal{A} = \Xs$, $\mathcal{Y} = \R$ and
$\kappa(x) = \Normal(\ip{x}{\theta},1)$),
\[
  \KL\big(\Prob_\theta^T\,\big\|\,\Prob_{\theta'}^T\big)
  = \frac12\sum_{t<T}\E_\theta\big[\ip{x_t}{\theta-\theta'}^2\big]
  \le \frac{T}{2}\Big(\sup_{x\in\Xs}\norm{x}\Big)^2\norm{\theta-\theta'}^2 < \infty .
\]
\end{corollary}

\begin{proof}
\leanok
\uses{thm:divergence_decomposition, lem:kl_gaussian}
The subtype $\Xs$ of $\R^d$ is countably generated, so Definition~\ref{def:env_pair}
applies with $\mathcal{A} = \Xs$ (compactness and spanning play no role there). The reward
kernels are $\kappa(x) = \Normal(\ip{x}{\theta},1)$ and
$\kappa'(x) = \Normal(\ip{x}{\theta'},1)$, so by Lemma~\ref{lem:kl_gaussian}
$\KL(\kappa(x)\|\kappa'(x)) = \tfrac12\ip{x}{\theta-\theta'}^2 \le \tfrac12 R^2\norm{\theta-\theta'}^2$
with $R := \sup_{x\in\Xs}\norm{x} < \infty$ (compactness). Apply
Theorem~\ref{thm:divergence_decomposition} in its finite form.
\end{proof}

\textbf{Lean remark.} The action type is the subtype $\Xs$ of \texttt{EuclideanSpace ℝ (Fin d)}
and the kernel is \texttt{fun x ↦ gaussianReal ⟪x, θ⟫ 1}, measurable by
\texttt{measurable\_gaussianReal} composed with the continuous inner product. The expectation
$\E_\theta[\ip{x_t}{\theta-\theta'}^2]$ is an integral of a bounded measurable function of the
history, so no integrability side condition arises.
